\section{Related Work}
\label{sec:related-ra}

The reversible-action substrate is bounded by five literatures:
transaction processing and the saga pattern, capability revocation in
operating systems, endpoint detection and response engineering, the
reversible-computing line, and the verified-systems pattern.

\paragraph{Compensating transactions and sagas.}
The saga pattern, introduced for long-running transactions whose
commit boundary is weaker than serializable schedules, pairs each
forward step with a compensating step that semantically undoes the
forward step~\cite{garciamolina1987sagas}. Workflow nets extend the model with
timed arcs and soundness criteria deciding whether every reachable
marking has a path to a final marking. The reversible-action substrate
inherits the saga vocabulary on the four reversible action variants:
the inverse executor is the compensating step, and the rollback
receipt is the typed artifact recording its completion. The substrate
does not inherit the workflow-net machinery; it has no timed
transition system beyond the monotonic counter, and the soundness
criterion the workflow-net literature proves is a strict
generalization the substrate does not claim. The contribution relative
to sagas is the constitutional-admission layer above the compensating
step: a rollback receipt is admissible only when the polity's
predicates accept its canonical body, which the saga literature treats
as a workflow-engine concern rather than as a polity-level decision.

\paragraph{Capability revocation in operating systems.}
EROS, seL4, and Capsicum treat capability revocation as a typed
operation on a capability token. The capability becomes unusable after
revocation, and the revocation event is a checkable artifact at the
machine boundary. The reversible-action substrate lifts the same
discipline to the cross-organizational boundary where the receipt
graph crosses kernels: a revoke-grant action is destructive, requires
bilateral cosignature, and the constitutional admission predicate
gates the revocation. The OS-capability literature solves revocation
at one kernel; the substrate gates revocation at the receipt-graph
boundary between kernels.

\paragraph{Endpoint detection and response.}
Vendor EDR engines have shipped TTL fields and manual rollback
workflows for a decade. CrowdStrike Falcon, SentinelOne, Microsoft
Defender for Endpoint, and Sublime Security each record an action
identifier, target host, action class, TTL duration, and rollback
handle for response actions; auditors inspect the rollback handle by
hand when something needs to be reverted~\cite{edrVendorSchemasPlaceholder}. The
typed four-receipt grammar (request, execution, rollback,
acknowledgement) is a normalization of existing vendor schemas; the
constitutional admission predicate on the rollback receipt is the
substrate-level addition. The substrate is not a new EDR; it is a
typed discipline the existing engines could in principle adopt.

\paragraph{Provenance and lineage.}
Why-provenance, where-provenance, and how-provenance distinctions in
the database literature~\cite{cheney2009provenance} are structurally related to
the receipt graph's lineage relations. A rollback receipt is a how-
provenance fragment: it records the inverse executor's signed
assertion that the OS operation completed, and the receipt graph
records the lineage from request to execution to rollback. The
substrate does not extend the provenance literature; it specializes
it to the cross-kernel-admission case the prior substrate
\cite{programmableSovereigntyAnonymous} already named.

\paragraph{Reversible computing.}
Reversible computing originates with Bennett's logical-reversibility
construction \cite{bennett1973logical}, which shows that a Turing
machine can be made information-preserving by carrying a history tape
so that every forward step has a unique inverse. The Pendulum line
realized the discipline in hardware
\cite{vieri1999pendulum}: a reversible-logic processor whose
instruction set is closed under inversion and whose thermodynamic cost
per operation is bounded by Landauer's principle rather than by the
clock. The Janus programming language
\cite{yokoyama2007janus} formalizes the same discipline at the
source-language level: every statement has a syntactic inverse and the
language admits an invertible self-interpreter. The reversibility in
this literature is thermodynamic-physical: each forward step is
recoverable because the runtime preserves the information needed to
undo it. The reversibility carried by a reversible-action substrate is
structural-policy: each forward step is admissible at the polity's
predicates only when paired with a typed rollback witness that
references the original execution receipt and the inverse executor's
signed completion. The two senses share the word and the structural
shape (paired forward and inverse, history preserved as data), but the
substrate makes no thermodynamic claim and inherits no Landauer bound.

\paragraph{Verified systems.}
The Lean-based industrial verification pattern is established by AWS
Cedar (Lean specification plus differential-tested Rust runtime) and
SampCert (Lean-verified differential-privacy primitives deployed
inside AWS Clean Rooms)~\cite{cutler2024cedar}. The older systems-
verification line is the foundational reference set. CompCert
\cite{leroy2009compcert} established that a realistic optimizing C
compiler can be verified end-to-end against a formal semantics; seL4
\cite{klein2009sel4} extended the discipline to an operating-system
kernel, with the abstract specification, the C implementation, and the
binary all verified against each other; IronFleet
\cite{hawblitzel2015ironfleet} extended it to a distributed
state-machine-replicated service. These verified-systems efforts
establish that formal verification at the artifact level is feasible
for production-scale code. The contribution of the prior substrate
\cite{programmableSovereigntyAnonymous} and of the reversible-action
extension is at the substrate-policy level rather than the artifact
level: the Lean theorems condition the admission predicates and the
amendment refinement obligations on which the runtime depends, not the
runtime's compiled code itself. The treaty intersection and amendment
refinement theorems extend the policy-verification pattern to
cross-kernel admission; the TTL-bounded amendment chain theorem added
here completes the shape on the response side.

\paragraph{Agentic-AI safety.}
The alignment-evaluation literature (METR, UK AI Safety Institute
Inspect, ARC Evals)~\cite{alignmentEvalHarnessesPlaceholder} bounds what a model is likely
to do during pre-deployment evaluation; the reversible-action substrate
gates whether a deployed agent's destructive action crosses the trust
boundary at runtime. The two disciplines compose rather than compete:
an evaluation result is an input predicate to the polity's
constitution, and the constitution's response to a destructive action
includes the bilateral-cosignature reduction. Constitutional AI
conditions outputs at training; the substrate gates outputs at
deployment.
